\section{静电场}

\subsection{分布、电势与能量}

静电场满足
\[
 \nabla\times\vect E=0,
 \qquad \nabla\cdot\vect E=\frac\rho{\varepsilon_0}.
\]
电荷分布对区域 $\Omega$ 给出 $Q(\Omega)=\int_\Omega\rho\dd V$；更一般地，它对测试函数 $f$ 给出线性泛函
\[
 \varrho(f)=\int_{\mathbb R^3}\rho(\vect r)f(\vect r)\dd V.
\]
点值泛函 $f\mapsto f(0)$ 和导数泛函 $f\mapsto f'(0)$ 都是分布。这样的语言可以无歧义地处理点电荷和多极子。

\begin{advdefinition}{点电荷与 Dirac $\delta$ 分布}\label{advA4:def:point-charge}
位于 $\vect r_0$、电荷量为 $q$ 的点电荷定义为
\[
 \rho(\vect r)=q\delta^{(3)}(\vect r-\vect r_0),
 \qquad
 \int q\delta^{(3)}(\vect r-\vect r_0)f(\vect r)\dd V=qf(\vect r_0).
\]
一维中 $\delta$ 不是普通函数，而由
$\int\delta(x-x_0)f(x)\dd x=f(x_0)$ 定义；其导数满足
$\int\delta'(x)f(x)\dd x=-f'(0)$。
\end{advdefinition}

Lorentz 核
\[
 L_\epsilon(x)=\frac1\pi\frac\epsilon{x^2+\epsilon^2}
\]
在 $\epsilon\downarrow0$ 时收敛到 $\delta(x)$：它非负、积分为1，且任意固定远离零点的区域上质量趋于零。分部积分还给出 $L_\epsilon'\to\delta'$。

\begin{advexamplebox}{Dirac $\delta$ 的四种逼近表示}\label{advA4:ex:delta-representations}
证明分布意义下
\[
 \delta(x)=\lim_{\epsilon\downarrow0}\frac1{\sqrt{2\pi}\epsilon}
 \e^{-x^2/(2\epsilon^2)}
 =\lim_{K\to\infty}\frac{\sin Kx}{\pi x},
\]
并说明周期 Fourier 核和连续 Fourier 积分如何给出
\[
 \delta(x)=\frac1{2\pi}\int_{-\infty}^{\infty}\e^{\ii kx}\dd k.
\]
\end{advexamplebox}
\begin{advsolution}
对 Gaussian 核令 $x=\epsilon u$，则对任意连续快速衰减的 $f$，
\[
 \int\frac{\e^{-x^2/(2\epsilon^2)}}{\sqrt{2\pi}\epsilon}f(x)\dd x
 =\int\frac{\e^{-u^2/2}}{\sqrt{2\pi}}f(\epsilon u)\dd u
 \longrightarrow f(0).
\]
对 sinc 核，用 Fourier 反演：
\[
 \int\frac{\sin Kx}{\pi x}f(x)\dd x
 =\frac1{2\pi}\int_{-K}^{K}\!\int f(x)\e^{\ii kx}\dd x\dd k
 \longrightarrow f(0).
\]
长度为 $2\pi$ 的周期区间上，Dirichlet 核
$D_N(x)=\sum_{n=-N}^{N}\e^{\ii nx}$ 满足
$(2\pi)^{-1}\int_{-\pi}^{\pi}D_N(x)f(x)\dd x\to f(0)$（在满足 Dirichlet 条件时取左右极限平均）。令周期趋于无穷，离散波数和变成积分，即得连续 Fourier 表示。四种表达都只在分布意义成立，不能把 $\delta(0)=\infty$ 当普通函数值参与代数运算。
\end{advsolution}

\subsection{Cartesian 多极展开}

当场点满足 $r$ 大于电荷分布尺度时，平移算符给出
\[
 \frac1{|\vect r-\vect r'|}
 =\sum_{n=0}^\infty\frac{(-1)^n}{n!}
 x'_{i_1}\cdots x'_{i_n}
 \partial_{i_1}\cdots\partial_{i_n}\frac1r.
\]
定义位置张量与导数张量
\[
 R_n(\vect r)=\underbrace{\vect r\cdots\vect r}_{n},
 \qquad
 T_n(\vect r)=(-1)^nr^{2n+1}
 \underbrace{\nabla\cdots\nabla}_{n}\frac1r.
\]
递推式
\[
 T_{n+1}=(2n+1)\vect r,T_n-r^2\nabla T_n
\]
给出
\[
 T_0=1,\qquad T_1=\vect r,\qquad
 T_2=3\vect r\vect r-r^2\mathsf1,
\]
\[
 (T_3)_{abc}=15x_ax_bx_c
 -3r^2(x_a\delta_{bc}+x_b\delta_{ca}+x_c\delta_{ab}).
\]
由于 $\nabla^2(1/r)=0$（$r\ne0$），$T_n$ 对任意两指标缩并均为零，是完全对称无迹张量。定义普通与无迹多极矩
\[
 Q_n=\int\rho(\vect r')R_n(\vect r')\dd V',
 \qquad
 D_n=\int\rho(\vect r')T_n(\vect r')\dd V'.
\]
电势可写成
\[
 \boxed{\varphi(\vect r)=\frac1{4\pi\varepsilon_0}
 \sum_{n=0}^\infty\frac{(D_n,R_n(\vect r))}{n!r^{2n+1}}}.
\]
前几阶为总电荷 $D_0=Q$、偶极矩
$\vect p=D_1=\int\rho\vect r\dd V$ 和四极矩
\[
 \mathsf D=D_2=\int\rho(3\vect r\vect r-r^2\mathsf1)\dd V.
\]
相应的偶极与四极电势为
\[
 \varphi_1=\frac{\vect p\cdot\vect r}{4\pi\varepsilon_0r^3},
 \qquad
 \varphi_2=\frac{\vect r\cdot\mathsf D\cdot\vect r}
 {8\pi\varepsilon_0r^5}.
\]
偶极可由相反点电荷作无穷小分离得到，分布写成
$\rho=-\vect p\cdot\nabla\delta^{(3)}(\vect r)$；四极则由两个相反偶极再分离得到，写成
$\rho=\tfrac12Q_{ij}\partial_i\partial_j\delta^{(3)}$。在光滑外势中的相互作用能由 Taylor 展开直接给出
\[
 U=Q\varphi_{\rm ext}(0)+\vect p\cdot\nabla\varphi_{\rm ext}(0)
 +\frac1{2!}Q_{ij}\partial_i\partial_j\varphi_{\rm ext}(0)+\cdots.
\]

\begin{advexamplebox}{对称无迹张量的独立分量数}\label{advA4:ex:stf-dimension}
证明三维 $n$ 阶完全对称无迹张量有 $2n+1$ 个独立分量。
\end{advexamplebox}
\begin{advsolution}
$n$ 阶完全对称张量只由指标1、2、3各出现多少次决定，因此
\[
 \dim S^n(\mathbb R^3)=\binom{n+2}{2}.
\]
取迹把 $S^n$ 映到 $S^{n-2}$。这个映射是满射：给定任意
$B\in S^{n-2}$，先取对称化的 $\delta\odot B$，再减去有限个更高重迹项即可构造迹恰为 $B$ 的张量。因此无迹子空间是迹映射的核，秩--零化度定理给出
\[
 \dim\operatorname{STF}^n
 =\binom{n+2}{2}-\binom n2
 =\boxed{2n+1}.
\]
$n=1$ 得3个偶极分量，$n=2$ 得5个四极分量，也与“对称矩阵6个分量减去一个迹条件”一致。
\end{advsolution}

\begin{advtheorem}{点电偶极子的完整分布场}\label{advA4:thm:dipole-distribution-field}
点偶极子在全空间的电场是
\[
 \boxed{\vect E(\vect r)=\frac1{4\pi\varepsilon_0}
 \frac{3(\vect p\cdot\vect r)\vect r-r^2\vect p}{r^5}
 -\frac{\vect p}{3\varepsilon_0}\delta^{(3)}(\vect r)}.
\]
\end{advtheorem}
\noindent\textbf{证明}\quad 分布恒等式
\[
 \partial_i\partial_j\frac1r
 =\frac{3x_ix_j-r^2\delta_{ij}}{r^5}
 -\frac{4\pi}{3}\delta_{ij}\delta^{(3)}(\vect r)
\]
可通过挖去半径 $\epsilon$ 小球后分部积分两次证明；小球面项在
$\epsilon\downarrow0$ 时留下各向同性张量
$-(4\pi/3)\delta_{ij}f(0)$。偶极势
$\varphi=-(4\pi\varepsilon_0)^{-1}p_j\partial_j(1/r)$，故
$E_i=(4\pi\varepsilon_0)^{-1}p_j\partial_i\partial_j(1/r)$，代入即得。$\delta$ 项保证分布意义下
$\nabla\cdot\vect E=-(\vect p\cdot\nabla\delta)/\varepsilon_0$。\hfill$\square$

\begin{advexamplebox}{电偶极子的合力与力矩}\label{advA4:ex:dipole-force-torque}
证明在尺度上缓变的外电场中
\[
 \vect F=(\vect p\cdot\nabla)\vect E,
 \qquad \vect M=\vect p\times\vect E.
\]
\end{advexamplebox}
\begin{advsolution}
把偶极子看成位于 $\vect r\pm\vect a/2$ 的 $\pm q$，其中 $\vect p=q\vect a$。合力
\[
 \vect F=q\vect E(\vect r+\vect a/2)-q\vect E(\vect r-\vect a/2)
 =(\vect p\cdot\nabla)\vect E+O(a^2).
\]
两力对中心的力矩为
\[
 \vect M=\frac{\vect a}{2}\times q\vect E
 +\left(-\frac{\vect a}{2}\right)\times[-q\vect E]
 =\vect p\times\vect E.
\]
也可从势能 $U=-\vect p\cdot\vect E$ 得到
$\vect F=-\nabla U$；静电场无旋时
$\nabla(\vect p\cdot\vect E)=(\vect p\cdot\nabla)\vect E$。力矩使 $\vect p$ 转向 $\vect E$，稳定方向为同向。
\end{advsolution}

\begin{advexamplebox}{任意取向偶极子与孤立中性导体球}\label{advA4:ex:dipole-grounded-sphere}
半径为 $R$、总电荷为零的孤立导体球外，距球心 $a>R$ 处放置偶极矩 $\vect p$，它与连心线夹角为 $\theta$。令 $\vect p=p_\parallel\vect e_r+\vect p_\perp$。求像系、势能、力和力矩，并与接地球比较。
\end{advexamplebox}
\begin{advsolution}
先构造接地球的像系：令 $d=a^2-R^2$，像点位于 $b=R^2/a$，其上有
\[
 q'=\frac{Rp_\parallel}{a^2},\qquad
 \vect p'=\left(\frac Ra\right)^3
 (p_\parallel\vect e_r-\vect p_\perp).
\]
偶极像的总电荷为零，但像点电荷 $q'$ 一般非零。
原题要求球不带电，故还须在球心增加
\[
 q_0=-q'=-\frac{Rp_\parallel}{a^2}.
\]
中心像只在球面上增加常势，不破坏等势边界，同时使感应总电荷为零。
接地与孤立中性是两种边界条件，不可互换。

像场在真实偶极子位置为
\[
 \vect E_{\rm im}=\frac1{4\pi\varepsilon_0}
 (A p_\parallel\vect e_r+B\vect p_\perp),
\]
其中
\[
 A=\frac R{d^2}+\frac{2R^3}{d^3}-\frac R{a^4},
 \qquad B=\frac{R^3}{d^3}.
\]
因感应电荷随源线性建立，势能为
\[
 \boxed{U=-\frac12\vect p\cdot\vect E_{\rm im}
 =-\frac{A p_\parallel^2+B p_\perp^2}{8\pi\varepsilon_0}}.
\]
相对接地球，中心补偿像使 $U$ 增加 $Rp_\parallel^2/(8\pi\varepsilon_0a^4)$。
固定偶极矩大小和空间方向求 $-\nabla U$，径向及横向分量分别为
\[
 F_r=-\frac1{8\pi\varepsilon_0}
 \left[\frac{4Ra p_\parallel^2}{d^3}
 +\frac{6R^3a(2p_\parallel^2+p_\perp^2)}{d^4}
 -\frac{4R p_\parallel^2}{a^5}\right],
\]
\[
 \vect F_\perp=\frac{(A-B)p_\parallel\vect p_\perp}
 {4\pi\varepsilon_0a},\qquad
 \boxed{\vect F=F_r\vect e_r+\vect F_\perp}.
\]
任意倾斜偶极子还受到非径向力，不能只交出 $F_r$。
力矩可直接由 $\vect M=\vect p\times\vect E_{\rm im}$ 求出。
若以增加 $\theta$ 为正，其分量为
\[
 \boxed{M_\theta=-\frac{p^2\sin2\theta}{8\pi\varepsilon_0}
 \left[\frac R{d^2}+\frac{R^3}{d^3}-\frac R{a^4}\right]}.
\]
因为方括号为正，径向的 $\theta=0,\pi$ 是稳定取向。
若题目改为接地球，则删去 $q_0$，即所有式中删去来自 $R/a^4$ 的项。
令 $a=R+h$、$R\to\infty$，两种边界的局部极限都为
\[
 U\longrightarrow-\frac{2p_n^2+p_t^2}{64\pi\varepsilon_0h^3},
\]
其中 $p_n,p_t$ 是相对于平面的法向、切向分量。
\end{advsolution}

\subsection{电介质}

极化强度 $\vect P$ 是单位体积的电偶极矩。若每单位体积有 $n$ 个相同偶极子，$\vect P=n\vect p$。

\begin{advtheorem}{极化电荷}\label{advA4:thm:polarization-charge}
分段光滑极化分布产生的电场等效于体束缚电荷与界面束缚电荷
\[
 \boxed{\rho_b=-\nabla\cdot\vect P},
 \qquad
 \boxed{\sigma_b=(\vect P_1-\vect P_2)\cdot\vect n},
\]
其中 $\vect n$ 从介质1指向介质2。介质外 $\vect P_2=0$ 时
$\sigma_b=\vect P\cdot\vect n$。
\end{advtheorem}
\noindent\textbf{证明}\quad 偶极势叠加为
\[
 \varphi=\frac1{4\pi\varepsilon_0}
 \int_V\vect P(\vect r')\cdot
 \frac{\vect r-\vect r'}{|\vect r-\vect r'|^3}\dd V'
 =\frac1{4\pi\varepsilon_0}\int_V
 \vect P\cdot\nabla'\frac1{|\vect r-\vect r'|}\dd V'.
\]
对 $\vect r'$ 分部积分，体项为
$-\int(\nabla'\cdot\vect P)/R$，边界项为
$\oint(\vect P\cdot\vect n)/R$。界面两侧法向相反，合并即得跳跃式。\hfill$\square$

定义
\[
 \vect D=\varepsilon_0\vect E+\vect P,
 \qquad \nabla\cdot\vect D=\rho_f.
\]
线性各向同性介质中
$\vect P=\varepsilon_0\chi_e\vect E$，故
$\vect D=\varepsilon\vect E$、
$\varepsilon=\varepsilon_0(1+\chi_e)$，场能密度为
$w=\tfrac12\vect D\cdot\vect E$。

\begin{advexamplebox}{位移极化的极化率}\label{advA4:ex:displacement-polarization}
把分子视为电荷 $\pm Q$ 的一维谐振子，相对位移势能为 $kx^2/2$，数密度为 $n$。求自由能系数 $\alpha$ 与电极化率 $\chi_e$。
\end{advexamplebox}
\begin{advsolution}
平衡时 $kx=QE$，单分子偶极矩
$p=Qx=Q^2E/k$。所以
\[
 P=np=\frac{nQ^2}{k}E.
\]
若自由能写成 $f=f_0+\alpha P^2/2-\vect P\cdot\vect E$，极值条件为 $\alpha P=E$，因此
\[
 \boxed{\alpha=\frac{k}{nQ^2}},
 \qquad
 \boxed{\chi_e=\frac{P}{\varepsilon_0E}
 =\frac{nQ^2}{\varepsilon_0k}}.
\]
$k$ 越大越难极化，$n,Q$ 越大越易极化，趋势正确。三维各向同性束缚若三个方向劲度相同，结果逐方向相同。
\end{advsolution}

\begin{advexamplebox}{取向极化与 Langevin 函数}\label{advA4:ex:orientation-polarization}
数密度为 $n$ 的永久偶极矩 $p$ 在温度 $T$ 下达到热平衡。求极化强度的精确式和弱场极化率。
\end{advexamplebox}
\begin{advsolution}
取 $\vect E$ 为极轴，单偶极能量 $U=-pE\cos\theta$。令
$x=pE/(k_BT)$，配分函数的角部分为
\[
 Z_1=2\pi\int_0^\pi\e^{x\cos\theta}\sin\theta\dd\theta
 =4\pi\frac{\sinh x}{x}.
\]
平均投影
\[
 \langle p\cos\theta\rangle
 =p\partial_x\ln Z_1
 =p\left(\coth x-\frac1x\right)
 =p\mathcal L(x).
\]
因此
\[
 \boxed{P=np\mathcal L\!\left(\frac{pE}{k_BT}\right)}.
\]
弱场 $\mathcal L(x)=x/3+O(x^3)$，故
\[
 P=\frac{np^2}{3k_BT}E,\qquad
 \boxed{\alpha=\frac{3k_BT}{np^2}},
 \quad
 \boxed{\chi_e=\frac{np^2}{3\varepsilon_0k_BT}}.
\]
强场 $x\to\infty$ 时 $P\to np$，正确体现饱和而不是无限线性增长。
\end{advsolution}

\begin{advtheorem}{Helmholtz 分解}\label{advA4:thm:helmholtz}
设 $\vect F$ 在区域 $V$ 内足够光滑。定义
\[
 \Phi(\vect r)=\frac1{4\pi}\int_V
 \frac{\nabla'\cdot\vect F(\vect r')}{|\vect r-\vect r'|}\dd V'
 -\frac1{4\pi}\oint_{\partial V}
 \frac{\vect F(\vect r')\cdot\dd\vect S'}{|\vect r-\vect r'|},
\]
\[
 \vect A(\vect r)=\frac1{4\pi}\int_V
 \frac{\nabla'\times\vect F(\vect r')}{|\vect r-\vect r'|}\dd V'
 -\frac1{4\pi}\oint_{\partial V}
 \frac{\dd\vect S'\times\vect F(\vect r')}{|\vect r-\vect r'|}.
\]
则
\[
 \boxed{\vect F=-\nabla\Phi+\nabla\times\vect A}.
\]
在全空间且场充分衰减时边界项消失；无旋场可写成梯度，无散场可写成旋度。
\end{advtheorem}
\noindent\textbf{证明}\quad 对 $G=1/(4\pi|\vect r-\vect r'|)$ 使用
$\nabla^2G=-\delta^{(3)}(\vect r-\vect r')$ 和矢量恒等式
\[
 \nabla\times\nabla\times(G\vect F')
 =\nabla[\nabla\cdot(G\vect F')]-\nabla^2(G\vect F').
\]
对 $\vect r'$ 分部积分，把导数从 $\vect F'$ 移到 $G$，体积分的边界项恰好产生定理中的两个曲面积分；其余体积分合并为
$-\int_V(\nabla^2G)\vect F'\dd V'=\vect F(\vect r)$。这也说明若不规定边界行为，势并不唯一。\hfill$\square$

全空间静电势与电场因而为
\[
 \varphi(\vect r)=\frac1{4\pi\varepsilon_0}
 \int\frac{\rho(\vect r')}{|\vect r-\vect r'|}\dd V',
\]
\[
 \vect E(\vect r)=\frac1{4\pi\varepsilon_0}
 \int\frac{\rho(\vect r')(\vect r-\vect r')}{|\vect r-\vect r'|^3}\dd V'.
\]
环路积分 $\oint\vect E\cdot\dd\vect l=0$，通量满足
$\oint\vect E\cdot\dd\vect S=Q_{\rm in}/\varepsilon_0$。点电荷穿过任意开曲面的通量为
$q\Omega/(4\pi\varepsilon_0)$，其中
$\Omega=\int_S\dd\vect S\cdot\vect e_r/r^2$ 是有向立体角。

\begin{advexamplebox}{圆锥立体角与球坐标立体角元}\label{advA4:ex:solid-angle}
半顶角为 $\alpha$ 的圆锥从顶点截取单位球面。求其立体角、$\alpha$ 改变时的立体角元，并推广到球坐标一般面元。
\end{advexamplebox}
\begin{advsolution}
单位球面上圆锥对应 $0\le\theta\le\alpha$，故
\[
 \boxed{\Omega=\int_0^{2\pi}\!\int_0^\alpha
 \sin\theta\dd\theta\dd\phi=2\pi(1-\cos\alpha)}.
\]
微分得
$\dd\Omega=2\pi\sin\alpha\dd\alpha$。一般球坐标面元在半径 $r$ 的球面上为
\[
 \dd S=r^2\sin\theta\dd\theta\dd\phi,
 \qquad
 \boxed{\dd\Omega=\frac{\dd S}{r^2}
 =\sin\theta\dd\theta\dd\phi}.
\]
若开曲面法向与视线夹角为 $\beta$，则
$\dd\Omega=(\cos\beta/r^2)\dd S$。$\alpha=\pi$ 时得到 $4\pi$，全空间极限正确。
\end{advsolution}

\begin{advexamplebox}{匀强场与共轴电偶极子的电场线}\label{advA4:ex:dipole-uniform-field-lines}
匀强电场 $E_0\vect e_z$ 中放置位于原点、偶极矩平行或反平行于 $\vect e_z$ 的点偶极子。求轴对称电场线方程。
\end{advexamplebox}
\begin{advsolution}
令 $s=+1$ 表示同向，$s=-1$ 表示反向，$p>0$。总电势为
\[
 \varphi=-E_0r\cos\theta+s\frac{p\cos\theta}{4\pi\varepsilon_0r^2}.
\]
轴对称无旋场可用通量函数 $\Psi$ 表示
\[
 E_r=\frac1{r^2\sin\theta}\partial_\theta\Psi,
 \qquad
 E_\theta=-\frac1{r\sin\theta}\partial_r\Psi.
\]
对匀强场和偶极场分别积分，得到
\[
 \Psi_{\rm uni}=\frac12E_0r^2\sin^2\theta,
 \qquad
 \Psi_{\rm dip}=s\frac{p}{4\pi\varepsilon_0r}\sin^2\theta.
\]
电场线沿线 $\Psi$ 不变，因此
\[
 \boxed{\sin^2\theta\left(
 \frac12E_0r^2+s\frac{p}{4\pi\varepsilon_0r}
 \right)=C}.
\]
同向时两条轴上偶极场都与匀强场同向，因而轴上无零点；赤道面有半径 $r=(p/4\pi\varepsilon_0E_0)^{1/3}$ 的零场环。反向时零点在正、负两条轴上，均满足 $r=(2p/4\pi\varepsilon_0E_0)^{1/3}$。把 $p$ 置零退回平行直线，把 $E_0$ 置零得到偶极场线 $r\propto\sin^2\theta$。
\end{advsolution}

静电能可写成
\[
 U=\frac12\int\rho\varphi\dd V.
\]
若 $\rho=\rho_1+\rho_2$、$\varphi=\varphi_1+\varphi_2$，交叉项是两个体系的相互作用能。

\begin{advtheorem}{Green 互易关系}\label{advA4:thm:green-reciprocity}
对在同一边界条件下的两组静电解，若边界项为零，则
\[
 \boxed{\int_V\rho_1\varphi_2\dd V=
 \int_V\rho_2\varphi_1\dd V}.
\]
\end{advtheorem}
\noindent\textbf{证明}\quad Green 第二恒等式给出
\[
 \int_V(\varphi_1\nabla^2\varphi_2-\varphi_2\nabla^2\varphi_1)\dd V
 =\oint_{\partial V}(\varphi_1\partial_n\varphi_2-\varphi_2\partial_n\varphi_1)\dd S.
\]
共同的齐次 Dirichlet 或 Neumann 边界使右边为零；代入
$\nabla^2\varphi_i=-\rho_i/\varepsilon_0$ 即得。\hfill$\square$

\begin{advtheorem}{静电场能}\label{advA4:thm:electrostatic-energy}
在场充分衰减或边界项消失时，
\[
 \boxed{\frac12\int\rho\varphi\dd V
 =\frac{\varepsilon_0}{2}\int E^2\dd V}.
\]
\end{advtheorem}
\noindent\textbf{证明}\quad 用 $\rho=-\varepsilon_0\nabla^2\varphi$ 分部积分：
\[
 \frac12\int\rho\varphi\dd V
 =-\frac{\varepsilon_0}{2}\oint\varphi\nabla\varphi\cdot\dd\vect S
 +\frac{\varepsilon_0}{2}\int|\nabla\varphi|^2\dd V.
\]
边界项消失且 $\vect E=-\nabla\varphi$，即得结论。场能非负，但点电荷自能发散；比较不同构型时应分离并扣除与位置无关的自能。

\begin{advexamplebox}{带电线段、旋转椭球与等势面}\label{advA4:ex:spheroid}
均匀带电线段位于 $z$ 轴 $-c<z'<c$。完成以下讨论：
\begin{enumerate}
 \item 证明等势面是长旋转椭球，电场线是正交的双曲线；
 \item 求长半轴 $A$、短半轴 $B$ 的孤立导体长旋转椭球电容；
 \item 写出其表面电荷密度；
 \item 给出扁旋转椭球的对应结果和等效像电荷分布；
 \item 在长轴上距中心 $l>A$ 处放点电荷 $Q$ 并把椭球接地，求椭球总感应电荷。
\end{enumerate}
\end{advexamplebox}
\begin{advsolution}
记到线段两端的距离为 $r_\pm$。线电荷密度为 $\lambda$ 时
\[
 \varphi=\frac\lambda{4\pi\varepsilon_0}
 \int_{-c}^{c}\frac{\dd z'}{\sqrt{\rho^2+(z-z')^2}}
 =\frac\lambda{4\pi\varepsilon_0}
 \ln\frac{r_++r_-+2c}{r_++r_--2c}.
\]
令长旋转椭球坐标
\[
 \xi=\frac{r_++r_-}{2c},\qquad
 \eta=\frac{r_+-r_-}{2c},
\]
则 $\varphi$ 只依赖 $\xi$；$\xi=$ 常数是焦点为 $\pm c$ 的长旋转椭球，$\eta=$ 常数是与它正交的双曲面，因此后者给出电场线。

导体表面 $\xi=\xi_0=A/c$，其中 $c=\sqrt{A^2-B^2}$。令
$Q_0(\xi)=\tfrac12\ln[(\xi+1)/(\xi-1)]$，外部势为
\[
 \varphi(\xi)=V\frac{Q_0(\xi)}{Q_0(\xi_0)}.
\]
远处 $Q_0(\xi)\sim1/\xi=c/r$，与 $Q/(4\pi\varepsilon_0r)$ 比较得到
\[
 \boxed{C_{\rm prolate}=\frac QV
 =\frac{4\pi\varepsilon_0c}{Q_0(A/c)}
 =\frac{8\pi\varepsilon_0c}
 {\ln[(A+c)/(A-c)]}}.
\]
用尺度因子
$h_\xi=c\sqrt{(\xi^2-\eta^2)/(\xi^2-1)}$，表面法向场为
\[
 E_n=-\frac1{h_\xi}\partial_\xi\varphi
 =\frac{V}{cQ_0(\xi_0)
 \sqrt{(\xi_0^2-1)(\xi_0^2-\eta^2)}}.
\]
故
\[
 \boxed{\sigma(\eta)=\varepsilon_0E_n}.
\]
它在尖端 $|\eta|=1$ 较大，符合曲率越大电荷越集中的物理图像。

扁旋转椭球取极半轴 $A$、赤道半轴 $B>A$，令
$c=\sqrt{B^2-A^2}$。把 $c$ 作解析延拓可得
\[
 \boxed{C_{\rm oblate}=
 \frac{4\pi\varepsilon_0c}{\arctan(c/A)}}.
\]
它可由焦盘 $s<c$ 上的等效电荷产生，盘面密度可取
\[
 \sigma_{\rm disk}(s)=\frac{Q}{2\pi c\sqrt{c^2-s^2}},
 \qquad 0<s<c,
\]
其积分为 $Q$，所产生的等势面正是共焦扁旋转椭球。$A\to B$ 时两种电容都趋于 $4\pi\varepsilon_0A$，球极限通过。

最后用互易定理。令状态1为外点电荷 $Q$ 与接地椭球，状态2为椭球保持单位电势且无外点电荷。状态2在轴上 $l$ 处的势为
\[
 \varphi_2(l)=\frac{Q_0(l/c)}{Q_0(A/c)}.
\]
互易关系给出 $Q\varphi_2(l)+Q_{\rm ind}\cdot1=0$，所以
\[
 \boxed{Q_{\rm ind}=-Q\frac{Q_0(l/c)}{Q_0(A/c)}
 =-Q\frac{\ln[(l+c)/(l-c)]}{\ln[(A+c)/(A-c)]}}.
\]
$l\to\infty$ 时感应总电荷趋零，$l\downarrow A$ 时趋于 $-Q$，极限合理。
\end{advsolution}

\subsection{静电边值问题与格林函数}

Poisson 方程与无源区的 Laplace 方程为
\[
 \nabla^2\varphi=-\frac\rho{\varepsilon_0},
 \qquad \nabla^2\varphi=0.
\]
Dirichlet 问题给定 $\varphi|_{\partial V}$，Neumann 问题给定
$\partial_n\varphi|_{\partial V}$；后者还须满足由 Gauss 定理得到的相容条件，并且解只确定到一个加法常数。

Green 函数满足
\[
 \nabla'^2G(\vect r,\vect r')=-4\pi\delta^{(3)}(\vect r-\vect r').
\]
Dirichlet Green 函数取 $G_D=0$ 于边界，解为
\[
 \varphi(\vect r)=\frac1{4\pi\varepsilon_0}\int_V
 \rho(\vect r')G_D(\vect r,\vect r')\dd V'
 -\frac1{4\pi}\oint_{\partial V}
 \varphi(\vect r')\partial_{n'}G_D\dd S'.
\]
Neumann Green 函数取常数法向导数以满足积分相容条件；选定势的加权平均后，解由体电荷项和
$G_N\partial_n\varphi$ 的边界项给出。

\begin{advtheorem}{静电 Green 函数的对称性}\label{advA4:thm:green-symmetry}
在相同的齐次边界条件和同一规范常数约定下，
\[
 \boxed{G_D(\vect r,\vect r')=G_D(\vect r',\vect r)},
 \qquad
 \boxed{G_N(\vect r,\vect r')=G_N(\vect r',\vect r)}.
\]
\end{advtheorem}
\noindent\textbf{证明}\quad 在 Green 第二恒等式中取
$u(\vect x)=G(\vect r,\vect x)$、
$v(\vect x)=G(\vect r',\vect x)$。体积分利用各自的 $\delta$ 源，等于
$-4\pi G(\vect r',\vect r)+4\pi G(\vect r,\vect r')$。Dirichlet 情形边界上两函数都为零；Neumann 情形两者法向导数相同为规定常数，边界项在固定平均值规范下也相消，所以体积分为零。\hfill$\square$

\begin{advexamplebox}{上半空间的 Dirichlet 与 Neumann Green 函数}\label{advA4:ex:half-space-green}
区域为 $z>0$，源点为 $\vect r_0=(x_0,y_0,z_0)$，镜像点为
$\vect r_0^*=(x_0,y_0,-z_0)$。求两类 Green 函数并解释像电荷法。
\end{advexamplebox}
\begin{advsolution}
令 $R_-=|\vect r-\vect r_0|$、$R_+=|\vect r-\vect r_0^*|$。在平面 $z=0$ 上 $R_-=R_+$，故
\[
 \boxed{G_D=\frac1{R_-}-\frac1{R_+}}
\]
满足 $G_D=0$。它等价于在镜像点放置反号像电荷，适用于接地平面。

对 Neumann 条件，平面上两个核的法向导数互为相反数，所以
\[
 \boxed{G_N=\frac1{R_-}+\frac1{R_+}}
\]
满足 $\partial_nG_N=0$，等价于同号像电荷。前者使势为奇延拓，后者使势为偶延拓。代回
$\nabla^2G=-4\pi\delta$ 时，镜像点在区域外，不引入额外真实源。Green 函数均关于 $\vect r,\vect r_0$ 对称。
\end{advsolution}

\begin{advexamplebox}{球外区域的两类 Green 函数}\label{advA4:ex:sphere-green}
对半径 $R$ 的球外区域 $r>R$，源点 $r_0>R$。求 Dirichlet Green 函数；给出零法向导数 Neumann Green 函数的球谐展开，并说明为什么后者不能由单个像点电荷完成。
\end{advexamplebox}
\begin{advsolution}
令
\[
 \vect r_0^*=\frac{R^2}{r_0^2}\vect r_0.
\]
在球面上有
$|\vect r-\vect r_0^*|=(R/r_0)|\vect r-\vect r_0|$，故
\[
 \boxed{G_D(\vect r,\vect r_0)=
 \frac1{|\vect r-\vect r_0|}
 -\frac{R/r_0}{|\vect r-\vect r_0^*|}}.
\]
第二项是位于球内的像电荷 $q^*=-(R/r_0)q$，因此不在求解区域内。

对 Neumann 条件，把轴取在 $\vect r_0$ 方向。自由核在球面附近展开为
\[
 \frac1{|\vect r-\vect r_0|}
 =\sum_{l=0}^\infty\frac{r^l}{r_0^{l+1}}P_l(\cos\gamma)
 \quad (r<r_0).
\]
加入在区域内无奇点、在无穷远衰减的项 $a_lr^{-l-1}P_l$，令
$\partial_rG_N|_{r=R}=0$，逐阶得到
\[
 a_l=\frac{l}{l+1}\frac{R^{2l+1}}{r_0^{l+1}},\qquad l\ge1.
\]
因此可取
\[
 \boxed{G_N=\frac1{|\vect r-\vect r_0|}
 +\sum_{l=1}^\infty\frac{l}{l+1}
 \frac{R^{2l+1}}{r^{l+1}r_0^{l+1}}P_l(\cos\gamma)}.
\]
单个像点的展开系数只能按固定几何比变化，不可能产生逐阶的 $l/(l+1)$，所以 Neumann 球问题通常需要连续像分布或无穷级数。
\end{advsolution}

\begin{advexamplebox}{两半球维持相反电势}\label{advA4:ex:split-sphere-potential}
半径 $R$ 的球面上，北半球维持 $+V$、南半球维持 $-V$，赤道绝缘缝宽忽略。求球内电势。
\end{advexamplebox}
\begin{advsolution}
边值只依赖 $\theta$，球内正则轴对称解为
\[
 \varphi(r,\theta)=\sum_{l=0}^\infty A_l
 \left(\frac rR\right)^lP_l(\cos\theta).
\]
令 $u=\cos\theta$，边值 $f(u)=V\operatorname{sgn}u$ 是奇函数，所以偶数 $l$ 系数为零。由 Legendre 正交性，对奇数 $l$，
\[
 A_l=\frac{2l+1}{2}\int_{-1}^1f(u)P_l(u)\dd u
 =V(2l+1)\int_0^1P_l(u)\dd u.
\]
利用
$\int P_l\dd u=[P_{l+1}-P_{l-1}]/(2l+1)$，得到
\[
 \boxed{A_l=V[P_{l-1}(0)-P_{l+1}(0)]\quad(l\text{ 奇})}.
\]
故
\[
 \boxed{\varphi=V\sum_{n=0}^\infty
 [P_{2n}(0)-P_{2n+2}(0)]
 \left(\frac rR\right)^{2n+1}P_{2n+1}(\cos\theta)}.
\]
前两项是
$\varphi=(3V/2)(r/R)P_1-(7V/8)(r/R)^3P_3+\cdots$。中心势为0，符合南北反对称；在任意 $r<R$ 级数绝对收敛。
\end{advsolution}

\begin{advexamplebox}{六面给定电势的立方体中心}\label{advA4:ex:cube-center-potential}
立方体两相对面维持 $+V$，其余四面维持 $-V$，棱边窄缝忽略。求中心电势。
\end{advexamplebox}
\begin{advsolution}
Dirichlet 解对边值线性。中心关于立方体的六个面完全对称，因此从任一面传到中心的调和测度相同；六个权重之和为1，所以每个面权重是 $1/6$。棱和顶点的集合面积为零，不影响调和测度。于是
\[
 \boxed{\varphi(0)=\frac16[2(+V)+4(-V)]=-\frac V3}.
\]
也可把六个“单独一面为1、其余面为0”的解相加：常数边界1的解处处为1，中心对称性迫使六个单面解在中心相等，故各为 $1/6$。这避免了不必要的三重 Fourier 级数。
\end{advsolution}

\begin{advtheorem}{球谐函数加法定理}\label{advA4:thm:spherical-addition}
单位球面上两方向 $(\theta,\phi)$ 与 $(\theta',\phi')$ 夹角为 $\gamma$，则
\[
 \boxed{P_l(\cos\gamma)=\frac{4\pi}{2l+1}
 \sum_{m=-l}^{l}Y_{lm}(\theta,\phi)Y_{lm}^*(\theta',\phi')}.
\]
\end{advtheorem}
\noindent\textbf{证明}\quad 固定第二方向，右边是关于第一方向的 $l$ 阶球谐函数，并在同时旋转两方向时不变，因此只能是 $C_lP_l(\cos\gamma)$。令两方向重合并在球面上积分：正交归一性给出右边积分为
$\sum_m1=2l+1$；而旋转对称性说明
$\sum_m|Y_{lm}|^2$ 是常数，其球面积分为 $2l+1$，故常数为 $(2l+1)/(4\pi)$。于是 $C_l=(2l+1)/(4\pi)$ 的倒数，得到公式。\hfill$\square$

\begin{advexamplebox}{接地球内的均匀带电直径}\label{advA4:ex:charged-diameter-in-sphere}
半径 $R$ 的接地导体球内，沿 $z$ 轴直径放置总电荷 $Q$ 的均匀细线。求球内电势的一种完整表示。
\end{advexamplebox}
\begin{advsolution}
线密度 $\lambda=Q/(2R)$。球内 Dirichlet Green 函数对位于
$\vect r'=z'\vect e_z$ 的源点为
\[
 G_D(\vect r,\vect r')=
 \frac1{|\vect r-\vect r'|}
 -\frac{R/|z'|}{|\vect r-(R^2/z')\vect e_z|},
\]
其中 $z'=0$ 时第二项取连续极限 $1/R$。因此
\[
 \boxed{\varphi(\vect r)=\frac{\lambda}{4\pi\varepsilon_0}
 \int_{-R}^{R}\left[
 \frac1{|\vect r-z'\vect e_z|}
 -\frac{R/|z'|}{|\vect r-(R^2/z')\vect e_z|}
 \right]\dd z'}.
\]
第二项表示球外轴线上的连续像电荷分布。对任意球面点
$|\vect r|=R$，反演恒等式
\[
 \left|\vect r-\frac{R^2}{z'}\vect e_z\right|
 =\frac R{|z'|}|\vect r-z'\vect e_z|
\]
使方括号逐点为零，所以边界条件被精确满足。积分关于 $z'$ 对称，故势关于赤道对称；线端点的理想化会造成局部奇异，但球内离开线段处势有限且满足 Laplace 方程。
\end{advsolution}

\subsection{静电平衡与电容矩阵}

设有若干导体。电势与各导体电荷量线性相关，可写成
$\varphi_i=P_{ij}q_j$。其逆关系为
$q_i=C_{ij}\varphi_j$，$\mathsf C=\mathsf P^{-1}$ 称电容矩阵；$C_{ii}$ 是自容，非对角项描述互容。

\begin{advtheorem}{静电能是电荷的正定二次型}\label{advA4:thm:capacitance-energy}
导体体系的静电能为
\[
 \boxed{U=\frac12P_{ij}q_iq_j
 =\frac12C_{ij}\varphi_i\varphi_j\ge0},
\]
且
\[
 \varphi_i=\frac{\partial U}{\partial q_i}.
\]
\end{advtheorem}
\noindent\textbf{证明}\quad 准静态地把电荷从零增加到最终值。在线性体系中，中间态 $sq_i$ 的电势为 $s\varphi_i$，所做功
\[
 U=\int_0^1(s\varphi_i)q_i\dd s=\frac12q_i\varphi_i.
\]
代入两种线性关系得到两个二次型。另一方面
$U=(\varepsilon_0/2)\int E^2\dd V\ge0$，除所有电荷为零外严格为正，所以 $\mathsf P$（在去掉整体规范自由度后）正定。对 $q_i$ 求偏导得到导体电势。\hfill$\square$

\begin{advtheorem}{电容矩阵的互易性}\label{advA4:thm:capacitance-reciprocity}
\[
 \boxed{P_{ij}=P_{ji},\qquad C_{ij}=C_{ji}}.
\]
\end{advtheorem}
\noindent\textbf{证明}\quad 取状态1仅第 $i$ 个导体带单位电荷，状态2仅第 $j$ 个导体带单位电荷。Green 互易关系给出状态1的第 $j$ 个导体电势等于状态2的第 $i$ 个导体电势，即
$P_{ji}=P_{ij}$。实对称可逆矩阵的逆仍对称，所以
$C_{ij}=C_{ji}$。等价地，二次型 $U$ 的混合二阶导数可交换。\hfill$\square$


\paragraph{参考电势、互容的符号与两导体模型。}
以上正定性采用无穷远电势为零并保留对无穷远的场能。若只保留两导体间的理想电容 $C$，则
\[
 \binom{q_1}{q_2}=C\begin{pmatrix}1&-1\\-1&1\end{pmatrix}
 \binom{\varphi_1}{\varphi_2},\qquad
 U=\frac C2(\varphi_1-\varphi_2)^2.
\]
此矩阵只有半正定性，共同平移电势不改变场，故不能直接求逆；选第二导体接地后，剩余关系才是 $q_1=C\varphi_1$。矩阵非对角元通常为负，而工程上把正数 $-C_{12}$ 称互电容，必须先说明约定。
以半径 $a<b$、长度 $L\gg b$ 的同轴圆柱为例，忽略端部，Gauss 定律给 $E(r)=q/(2\pi\varepsilon Lr)$，从而
\[
 \varphi_1-\varphi_2=\frac q{2\pi\varepsilon L}\ln\frac ba,
 \qquad C=\frac{2\pi\varepsilon L}{\ln(b/a)}.
\]
将场能积分 $\int_a^b\varepsilon E^2(2\pi rL)\dd r/2$ 代入，恰得 $q^2/(2C)$；窄隙 $b-a\ll a$ 时又趋于平行板结果 $\varepsilon(2\pi aL)/(b-a)$。
